\chapter{procint: Process Intelligence as Process Law}
\label{ch:procint}
\label{sec:procint}

Chapter~\ref{ch:mfact} developed \textsf{mfact} as a content-neutral
admission apparatus: candidates, typed refusals, admissions with covers
proofs, manifests, and the closed family of valid objections say nothing,
by themselves, about what is being admitted. This chapter turns to the
first domain manufactured through that apparatus: \textsf{procint}, a
Lean~4 formalization of process-intelligence semantics --- Petri nets,
workflow nets and their soundness property, alignment-based conformance
checking, and OCEL~2.0 object-centric event logs. \textsf{procint} is not
a demonstration corpus assembled to exercise \textsf{mfact}; it is the
substrate the rest of this thesis's empirical claims are about, and its
declaration catalog, drafted across independent LLM-assisted lanes and
admitted only through kernel testing, is exactly the kind of
manufactured-not-authored content the mfact gate exists to police. As
the accompanying paper's ``procint: Process Intelligence as Process Law''
section puts it, process
intelligence in this sense does not mean dashboarding over event logs; it
means a mathematical substrate for deciding, replaying, and explaining
process standing --- what happened, what was admitted, what was refused,
which object relations were involved, which conformance obligations
passed, and which receipts authorize the resulting claim.

\textsf{procint} depends on Mathlib~\cite{mathlib2020} for its background
mathematics and on CSLib~\cite{cslib2026} for transition-system and trace
infrastructure, and is organized into namespaces that separate concerns
cleanly along the process-intelligence canon: \texttt{ProcInt.Foundations}
for metrics, identifiers, and labeled transition systems;
\texttt{ProcInt.Petri} for nets, firing, reachability, the state
equation, invariants, boundedness, liveness, object-centric Petri nets,
stochastic nets, and unfoldings; \texttt{ProcInt.Workflow} for workflow
nets, the short-circuit construction, and soundness; \texttt{ProcInt.Logs}
for events, traces with monotone timestamps, event logs, and XES;
\texttt{ProcInt.Models} for process trees, POWL, directly-follows graphs,
causal nets, BPMN, and Declare; \texttt{ProcInt.Conformance} for moves,
alignments, cost, fitness, and token replay; \texttt{ProcInt.Ocel} for
OCEL~2.0 events, objects, event-to-object and object-to-object relations,
flattening, and lifecycles; \texttt{ProcInt.Ocpq} for object-centric
process queries; \texttt{ProcInt.Analytics} for temporal, causal,
correlational, cube, and streaming analytics; and \texttt{ProcInt.Registry}
for the cross-product algorithm and model-breed registries that are
themselves data-driven rows of the same declaration graph rather than
hand-written tables. This chapter takes the workflow-net soundness
development as its center of gravity, because it is where the corpus's
single crown-jewel theorem lives, but treats the surrounding layers ---
Petri semantics beneath it, conformance and OCEL~2.0 beside it --- with
enough precision to support the canon-coverage claims the accompanying
paper's canon-status table makes.

\section{Petri Net Semantics}
\label{sec:procint-petri}

The base layer of \textsf{procint}'s process semantics follows Murata's
survey treatment of Petri nets~\cite{murata1989}, formalized directly
rather than via an intermediate encoding. A net is places, transitions,
and weighted flow over finite types, and a marking is a finitely
supported function \(M : P \to_{f} \mathbb{N}\) rather than a bare
function into the naturals; representing markings as \texttt{Finsupp}
values, Mathlib's finitely supported function type, does two things at
once. First, it unifies the weighted-net and colored-net cases: a
weighted marking is a \texttt{Finsupp} into \(\mathbb{N}\), and richer
codomains extend the same representation without a second definitional
layer. Second, it gives the whole marking algebra --- firing, the state
equation \(M' = M + A^{\mathsf T}u\) relating a reached marking to the
initial marking via the incidence matrix \(A\) and a firing-count vector
\(u\), and place/transition invariants --- as lemmas already available
over \texttt{Finsupp} in Mathlib, with truncated subtraction standing in
for the natural-number semantics of ``consume what is present.'' This
representational choice recurs later in this chapter: the same
truncated-subtraction discipline is what the crown-jewel theorem's
countermodel exploits, and what the correspondence rail of
Chapter~\ref{ch:correspondence} has to reconcile against Rust's own
saturating arithmetic.

Firing, reachability, boundedness, and liveness are defined over this
marking algebra in \texttt{ProcInt.Petri} and reused, rather than
redefined, by the workflow-net layer above it and by the object-centric
Petri net (OCPN) and stochastic-net extensions that follow
the accompanying paper's canon-status table's ``formalized''/``seed
instance'' distinction: the
core marking, firing, and invariant theory is formalized in full, while
OCPN and stochastic nets carry seed instances built on top of it rather
than independent full developments.

\section{Workflow Nets and the Soundness Property}
\label{sec:procint-wfnet}

A \texttt{WfNet} carries a Petri net as a field --- composition, not
inheritance by extension --- together with the source-place, sink-place,
and connectivity conditions of van der Aalst's Definition~6--7
\cite{vanderaalst1997}: a single source place with no input transitions, a
single sink place with no output transitions, and every node on some path
from source to sink. This composition-over-extension choice keeps the
underlying Petri net's marking algebra directly reusable rather than
re-derived under a workflow-specific wrapper.

The classical soundness property of a workflow net is usually stated as
one conjunction: from the initial marking, the final marking is always
reachable (option to complete); when the final marking is reached, no
tokens remain elsewhere (proper completion); and every transition can
still fire from some reachable marking (no dead transitions). Following
the observation that these three clauses are logically independent of one
another, \textsf{procint} does not encode soundness as a single
conjunction-shaped definition. Instead it states option-to-complete,
proper-completion, and no-dead-transitions as three separately named
predicates, with \texttt{Sound} defined as their explicit conjunction.
This keeps partial results honest in a literal, checkable sense: a lemma
that establishes one clause is stated, and can be cited, about exactly
that clause, and the fact that \S\ref{sec:procint-crown} below can report
\texttt{PROVEN\_SUPPORT} status for two of the three clauses separately
from the equivalence that uses all three depends directly on this
decomposition.

\section{The Crown-Jewel Theorem}
\label{sec:procint-crown}

The single most load-bearing result in the \textsf{procint} corpus is van
der Aalst's soundness characterization, Lemma~8 of
\cite{vanderaalst1997}: for a workflow net satisfying the source/sink/
connectivity conditions above, \texttt{Sound} holds of the net if and only
if the short-circuited net --- the net obtained by adding a transition
from the sink place back to the source place, folding the ``process
completes and a new instance may start'' behavior into an ordinary
liveness-and-boundedness question --- is both live (every transition can
recur from every reachable marking) and bounded (the reachable marking
set is finite). \textsf{procint} calls this the crown-jewel theorem
because the release's Standing Quadrature and manifest apparatus
(Chapter~\ref{ch:empirical}) is organized around it as the single result
whose kernel-admitted status the release most directly advertises, and
because it is the result that most clearly demonstrates the pipeline
producing a theorem whose informal target predates the pipeline by three
decades, rather than a theorem invented to be easy to formalize.

\begin{theorem}[Crown-jewel soundness equivalence, informal statement]
\label{thm:crown-informal}
Let \(W\) be a workflow net over a finite transition type \(T\) satisfying
the source, sink, and on-path connectivity conditions of van der Aalst's
Definition~6--7. Then \(W\) is sound if and only if the short-circuited
net \(W^{sc}\) is live and bounded.
\end{theorem}

The obligation ladder recorded in the repository's crown research
tracking (\texttt{research/wfnet/obligations.toml}) states the run's
standing toward this theorem as four ordered obligations. Order~0 is the
equivalence itself, annotated \emph{``the crown equivalence itself (van
der Aalst 1997, Lemma~8)''} and carrying status \texttt{PROVEN}. Orders~1
and~2, \emph{``soundness implies the final marking is reachable''} and
\emph{``soundness implies every transition can fire,''} carry status
\texttt{PROVEN\_SUPPORT}: they are the proper-completion and
no-dead-transitions directions of the soundness decomposition from
\S\ref{sec:procint-wfnet}, each proven separately and feeding the
equivalence rather than duplicating it. Order~3, discussed on its own in
\S\ref{sec:procint-open} below, is the one obligation in this ladder ---
and, as \S\ref{sec:procint-open} details, in the entire declaration
catalog --- that remains \texttt{STATED} rather than \texttt{PROVEN}. The
\texttt{[crown]} block of the same file records the release-level
judgment, \texttt{status = "PROVEN"}, computed the same way every other
status claim in this thesis is computed: by grepping the admitted TTL
declaration catalog for the relevant declarations' recorded status, never
by hand-asserting the release's headline claim. The generated fragment
\texttt{crown\_jewel\_status.tex} (quoted in full in
Chapter~\ref{ch:ai-development})
reports this same ground-truth classification directly from
\texttt{ProcInt.crownJewel\_status}, so this chapter's prose claim and the
release's generated claim are, structurally, the same read of the same
underlying value rather than two independently maintained assertions that
could drift apart.

\subsection{The \texttt{[Finite T]} Hypothesis and Its Countermodel}
\label{sec:procint-finite-t}

Theorem~\ref{thm:crown-informal} as stated above is not unconditionally
true of van der Aalst's original formulation once the transition set is
allowed to be infinite, and \textsf{procint}'s Lean statement of the
crown-jewel theorem is correspondingly narrower than the informal
statement: it carries an explicit \texttt{[Finite T]} instance hypothesis
on the transition type. This hypothesis is not a formalization
convenience adopted for proof-engineering ease; it is necessary, and the
repository carries a Lean-admitted witness proving that it is necessary
rather than merely asserting so in prose.

The witness is a purpose-built counterexample net, assembled in five
stages across the \texttt{Workflow.Countermodel} module (declared in
\texttt{packs/lean-math-pack/fragments/workflow\_countermodel.ttl}, module
order 34) and structured around a place type
\texttt{CrownCounterPlace} (an input place \texttt{i}, a queue place
\texttt{q}, an output place \texttt{o}, and a countably infinite family of
counter places \texttt{c n}) together with a transition type
\texttt{CrownCounterTransition := ℕ ⊕ ℕ}, one summand for an
\emph{absorb} transition at each index \(n\) and one for a matching
\emph{emit} transition. Absorb transition \(n\) consumes one token from
\texttt{i} and produces \(n\) tokens into \texttt{q} together with one
token into \texttt{c n}; emit transition \(n\) consumes those \(n\) tokens
from \texttt{q} and the one token from \texttt{c n}, and produces one
token into \texttt{o}. Because \(\mathbb{N} \oplus \mathbb{N}\) is
infinite --- witnessed directly by the instance declaration
\texttt{infinite\_CrownCounterTransition} --- this net has infinitely many
transitions, one absorb/emit pair per natural number, and can therefore
never satisfy \texttt{[Finite T]}.

The construction is engineered so that this infinite net is nonetheless
sound while its short-circuit is unbounded, the exact combination the
crown-jewel equivalence rules out under \texttt{[Finite T]}. Five
declarations carry this argument: \texttt{crownCounter\_reaches\_mid} and
\texttt{crownCounter\_reaches\_final} establish the reachability facts the
soundness argument needs; \texttt{crownCounter\_sound} establishes that
the net satisfies all three soundness clauses of
\S\ref{sec:procint-wfnet}; \texttt{crownCounter\_not\_bounded} establishes
that for every bound \(k\), some reachable marking of the short-circuit
violates it, because absorbing at index \(n\) drives \(n\) tokens into
\texttt{q} for arbitrarily large \(n\); and the main theorem,
\texttt{WfNet.infinite\_transition\_countermodel\_sound\_not\_bounded},
assembles these into the conjunction \(\mathrm{Infinite}\,
\text{CrownCounterTransition} \land \exists\, W,\, W.\mathrm{Sound}
\land \lnot \exists\, k,\, W.\mathrm{shortCircuit}.\mathrm{Bounded}\,
W.\mathrm{initialMarking}\,k\).

These five declarations are recorded in the catalog with \texttt{status
"stated"}, and this is a deliberate, permanent classification rather than
an unfinished proof obligation awaiting a future lane. The countermodel's
job is to witness a negative fact --- that the equivalence fails outside
\texttt{[Finite T]} --- and the repository's guardrails treat promoting
these five declarations to a proven, axiom-audited status as a specific
failure mode to be mechanically blocked rather than merely discouraged:
a countermodel family that quietly acquired \texttt{PROVEN} status would
misrepresent what the family is for. The corresponding refusal vocabulary
--- \texttt{WFNET\_INFINITE\_TRANSITION\_COUNTERMODEL} as the status key
for this theorem family, \texttt{countermodel\_not\_promoted} as the guard
name, and \texttt{COUNTERMODEL\_PROMOTION\_REFUSED} as the refusal fired
if a promotion is attempted --- exists precisely so that this
countermodel's status cannot silently drift from ``deliberate negative
witness'' to ``forgotten stub.'' Read correctly, then, these five
declarations are not a gap in the crown-jewel development; they are part
of its statement, in the same sense that a theorem's hypotheses are part
of its statement. \texttt{[Finite T]} is not a hedge adopted for proof
convenience --- it is exactly the hypothesis the countermodel shows is
required, made visible in the formal statement rather than left implicit
in prose the way an informal restatement of Lemma~8 might leave it.

\section{Conformance Checking and Alignments}
\label{sec:procint-conformance}

\textsf{procint}'s conformance-checking layer, \texttt{ProcInt.Conformance},
follows Carmona et al.'s alignment-based treatment of conformance checking
\cite{carmona2018} in preference to naive token-replay diagnostics.
Alignments relate an observed trace to a model run via four move kinds:
synchronous moves where the log and the model agree, log-only moves,
model-only moves, and silent (\(\tau\)) moves; a cost function over moves
induces optimal alignments as the alignments of least total cost. Fitness
--- the conformance metric this thesis returns to in
Chapter~\ref{ch:correspondence}, where it is the quantity a Rust-extracted
implementation and a hand-authored Lean model are related over --- is
represented as a \texttt{Between01} value by construction, so that the
boundedness of the metric (fitness always lies in \([0,1]\)) is a
type-level fact enforced by the representation itself rather than a
theorem that has to be separately proven and could, in principle, have
been forgotten.

\section{OCEL 2.0 Coverage}
\label{sec:procint-ocel}

Object-centric event logs follow the OCEL~2.0 specification
\cite{ocel2}, Definitions~1--2, in \texttt{ProcInt.Ocel}: typed events and
typed objects, event-to-object and object-to-object relations carrying
qualifiers, and attribute maps over both events and objects. Flattening
onto a classical, single-case-notion event log, and object lifecycle
extraction, are represented as derived operations over this base
representation rather than as independently axiomatized structures, so
that a flattening result is a theorem about the OCEL~2.0 representation
rather than an assumption stated separately about a flattened log.
Object-centric Petri nets, the net-level counterpart of object-centric
event data, follow~\cite{ocpn2020} and are formalized, per
the accompanying paper's canon-status table, as a seed instance built
on the Petri semantics of \S\ref{sec:procint-petri} rather than as a full
independent development at the same depth as the workflow-net soundness
theory.

\section{The One Open Obligation: Unfolding Correctness}
\label{sec:procint-open}

Of the \TotalDecls\ declarations recorded in the catalog at the time of
the release described by \texttt{release\_macros.tex},
exactly one is genuinely open proof
work rather than a completed theorem, a definition, a test oracle, or a
deliberately stated negative witness of the kind discussed in
\S\ref{sec:procint-finite-t}: order-3 of the crown obligation ladder,
\texttt{ProcInt.BranchingProcess.isUnfoldingOf\_statement}, tracked in
\texttt{research/wfnet/obligations.toml} under the name
\texttt{unfolding\_correctness} with status \texttt{STATED} and annotated
\emph{``branching-process unfolding statement (separate stated lane).''}
This is worth stating plainly rather than folding into the crown-jewel
narrative above, because it is a different kind of incompleteness from
the countermodel family: the countermodel is deliberately, permanently
stated as a negative witness that should never be promoted; the unfolding
statement is stated because the corresponding proof has not yet been
produced, and promoting it to \texttt{proven} remains an open task for a
future lane, not a status this repository's guardrails forbid in
principle.

The declaration, recorded in module \texttt{Petri.Unfolding} in
\texttt{packs/lean-math-pack/fragments/petriext.ttl}, states --- but does
not yet prove --- what it means for a branching process to be a legal
unfolding of a Petri net \(N\): given a branching process \(B\) over
places \(P\), transitions \(T\), condition labels \(C\), and events \(E\),
\(B\) is an unfolding of \(N\) exactly when \(B\)'s occurrence net is
acyclic, has no backward conflict, and its event/condition labeling
respects \(N\)'s pre- and post-set structure on singleton arcs:
\begin{center}
\begin{tabular}{l}
\(B.\mathrm{occ.Acyclic} \;\land\; B.\mathrm{occ.NoBackwardConflict}\)\\
\(\land\; (\forall\, e\, c,\; c \in B.\mathrm{occ.preC}\,e \to N.\mathrm{pre}\,(B.\mathrm{eventTrans}\,e)\,(B.\mathrm{condPlace}\,c) \neq 0)\)\\
\(\land\; (\forall\, e\, c,\; c \in B.\mathrm{occ.postC}\,e \to N.\mathrm{post}\,(B.\mathrm{eventTrans}\,e)\,(B.\mathrm{condPlace}\,c) \neq 0)\)
\end{tabular}
\end{center}
The declaration's own doc comment attributes the homomorphism condition
this statement specializes --- to ordinary nets, on singleton arcs --- to
Engelfriet's Definition~3.1 characterization of branching-process
unfoldings (Engelfriet 1991), as the source fragment states it directly.
No bibliographic entry for this attribution exists in this thesis's
reference list, and none is added here: inventing a bibtex record for a
citation that appears only as an in-source doc-comment attribution, absent
independent verification of its bibliographic details, would be exactly
the kind of fabrication this thesis's own manufacturing discipline is
designed to refuse elsewhere in the pipeline, and prose is not exempt
from that discipline merely because it is prose rather than a TTL
declaration. The attribution is therefore reported here as the source
states it --- an in-repository doc comment, not a peer-reviewed citation
this thesis independently verifies --- and readers who need the primary
source should consult the branching-process unfolding literature
directly rather than treat the name and year given here as a full
bibliographic record.

Because \texttt{isUnfoldingOf\_statement} is a \texttt{def} of type
\texttt{Prop} rather than a \texttt{theorem} with a completed proof term,
it participates in the corpus the way any stated-but-undischarged
statement does: it is visible in the declaration catalog, it is counted
among the \StatedCount\ stated declarations reported in
Table~\ref{tab:corpus} of \texttt{release\_macros.tex}, and it is excluded from the
\KernelTheorems\ kernel-admitted, axiom-audited theorems the release
advertises as proven. Its presence in the catalog at \texttt{STATED}
status, rather than its silent omission, is itself part of what this
thesis means by an honest, receipted release: a formalization effort that
reports only its completed theorems and omits its open statements would
overstate its own standing, and the manifest-and-quadrature apparatus of
Chapter~\ref{ch:empirical} is built specifically to make that kind of
omission structurally visible rather than merely discouraged in review.
